% Part: first-order-logic
% Chapter: models-theories
% Section: set-theory

\documentclass[../../../include/open-logic-section]{subfiles}

\begin{document}

% SEGMENT OLP-0172-S01

\olfileid{fol}{mat}{set}

\olsection{கணங்களின் கோட்பாடு}

கணங்களின் கோட்பாட்டில் கிட்டத்தட்ட எல்லாக் கணிதத்தையும் வளர்க்கலாம்.
இக்கோட்பாட்டில் கணிதத்தை வளர்ப்பதில் பல பணிகள் அடங்குகின்றன. முதலில்,
தொடர்பு~$\in$ க்கான அடிகோள்களின் தொகுப்பு தேவை. சிலவேளைகளில்~$\in$
க்கு முரண்படும் பண்புகளை வழங்கும் பல அடிகோள் அமைப்புகள் உருவாக்கப்பட்டுள்ளன.
தேர்வு அடிகோளுடன் கூடிய செர்மெலோ--ஃபிராங்கெல் கணக் கோட்பாடான
$\Log{ZFC}$ என்ற அடிகோள் அமைப்பு தனித்து நிற்கிறது. கணிதவியலாளர்கள்
நிரூபிக்க எதிர்பார்க்கும் ஏறத்தாழ அனைத்தையும் நிரூபிக்க அதன் அடிகோள்கள்
போதும் என்பதால், அதுவே மிகப் பரவலாகப் பயன்படுத்தப்பட்டும்
ஆராயப்பட்டும் வருகிறது. ஆனால் இதை நிறுவுவதற்கு முன், கணிதவியலாளர்கள்
கூற விரும்பும் அனைத்தையும் எவ்வாறு \emph{வெளிப்படுத்த} முடியும் என்பதைத்
தெளிவாக்க வேண்டும். தொடக்கமாக, மொழியில் !!{constant}s அல்லது
!!{function}s இல்லை. ஆகவே குறிப்பிட்ட கணங்களைப் (எடுத்துக்காட்டாக,
$\emptyset$ அல்லது $\Nat$) பற்றியோ, கணங்கள் மீதான செயல்களைப்
(எடுத்துக்காட்டாக, $X \cup Y$ மற்றும் $\Pow{X}$) பற்றியோ, தொடர்புகள்
மற்றும் சார்புகள் போன்ற கணங்கள் அல்லாதவற்றை உள்ளடக்கிய பிற
கட்டுமானங்களைப் பற்றியோ பேச முடியுமா என்பது முதல் பார்வையில் தெளிவில்லை.

முதலில், ``ஓர் உறுப்பு'' என்பது மட்டுமே நமக்கு வேண்டிய தொடர்பல்ல;
``ஓர் உட்கணம்'' என்பதும் ஏறத்தாழ அதே அளவு முக்கியமானது. ஆனால்
``ஓர் உறுப்பு'' என்பதைக் கொண்டு ``ஓர் உட்கணம்'' என்பதை
\emph{வரையறுக்க} முடியும். இதற்காக, கணக் கோட்பாட்டின் மொழியில்
!!a{formula}~$!A(x, y)$ ஒன்றைக் கண்டறிய வேண்டும்; கணங்களின் சோடி
~$\tuple{X, Y}$ அதை நிறைவுறுத்துவது iff $X \subseteq Y$. $X$ இன்
எல்லா !!{element}s உம்~$Y$ இன் !!{element}s ஆக இருந்தால் மட்டுமே
$X$ ஆனது $Y$ இன் உட்கணம். எனவே பின்வரும் வாய்பாட்டால் $\subseteq$ ஐ
வரையறுக்கலாம்:
\[
\lforall[z][(z \in x \lif z \in y)]
\]
இப்போது, ஒரு வாய்பாட்டில் தொடர்பு~$\subseteq$ ஐப் பயன்படுத்த
விரும்பும்போதெல்லாம், அதற்குப் பதிலாக அந்த வாய்பாட்டையே பயன்படுத்தலாம்
($x$ மற்றும் $y$ ஐ ஏற்றவாறு மாற்றியும், தேவைப்பட்டால் கட்டுண்ட
மாறி~$z$ ஐ மறுபெயரிட்டும்). எடுத்துக்காட்டாக, கணங்களின் விரிவுநிலை
என்பது எந்தக் கணங்கள்~$x$ மற்றும் $y$ ஒன்றையொன்று கொண்டிருந்தாலும்,
$x$ மற்றும் $y$ ஒரே கணமாக இருக்க வேண்டும் என்பதாகும். இதை
$\lforall[x][\lforall[y][((x \subseteq y \land y
    \subseteq x) \lif x = y)]]$ என்று வெளிப்படுத்தலாம். அல்லது மேலுள்ள
வரையறையால் $\subseteq$ ஐ மாற்றினால், பின்வருமாறு எழுதலாம்:
\[
\lforall[x][\lforall[y][((\lforall[z][(z \in x \lif z \in y)] \land
    \lforall[z][(z \in y \lif z \in x)]) \lif x = y)]].
\]
உண்மையில் இது $\Log{ZFC}$ இன் அடிகோள்களில் ஒன்றான
``விரிவுநிலை அடிகோள்''.

$\emptyset$ க்கான !!{constant} எதுவுமில்லை; ஆனால் ``$x$ வெறுமையானது''
என்பதை $\lnot \lexists[y][y \in x]$ என்று வெளிப்படுத்தலாம். அப்போது
``$\emptyset$ உள்ளது'' என்பது !!{sentence}~$\lexists[x][\lnot
\lexists[y][y \in x]]$ ஆகிறது. இது~$\Log{ZFC}$ இன் மற்றோர் அடிகோள்.
(ஒரே ஒரு வெற்றுக் கணம் மட்டுமே உள்ளது என்பதை விரிவுநிலை அடிகோள்
உட்கிடையச் செய்கிறது.) கணக் கோட்பாட்டின் மொழியில் $\emptyset$ ஐப்
பற்றிப் பேச விரும்பும்போதெல்லாம், ``வெறுமையான ஒரு கணம் உள்ளது,
மேலும் \dots'' என்று எழுதுவோம். எடுத்துக்காட்டாக, $\emptyset$ ஒவ்வொரு
கணத்தின் உட்கணம் என்பதை வெளிப்படுத்த, பின்வருமாறு எழுதலாம்:
\[
\lexists[x][(\lnot\lexists[y][y \in x] \land \lforall[z][x \subseteq
    z])]
\]
இங்கு $x \subseteq z$ என்பதையும் அதன் வரையறையால் மாற்ற வேண்டும்.

$X \cup Y$ மற்றும் $\Pow{X}$ போன்ற கணங்கள் மீதான செயல்களைப் பற்றிப்
பேச, இதேபோன்ற உத்தியைப் பயன்படுத்த வேண்டும். கணக் கோட்பாட்டின் மொழியில்
சார்புக் குறியீடுகள் இல்லை. ஆனால் சார்புத் தொடர்புகளான $X \cup Y = Z$
மற்றும் $\Pow{X} = Y$ ஆகியவற்றைப் பின்வருமாறு வெளிப்படுத்தலாம்:
\begin{align*}
& \lforall[u][((u \in x \lor u \in y) \liff u \in z)]\\
& \lforall[u][(u \subseteq x \liff u \in y )]
\end{align*}
$X \cup Y$ இன் !!{element}s என்பவை~$X$ இன் !!{element}s அல்லது~$Y$ இன்
!!{element}s ஆக இருக்கும் கணங்களே; மேலும் $\Pow{X}$ இன் !!{element}s
என்பவை~$X$ இன் உட்கணங்களே. இருப்பினும், $x \cup y$ அல்லது $\Pow{x}$
ஆகியவற்றைச் சொற்களைப் போலப் பயன்படுத்த இதனால் முடியாது; $X \cup Y = Z$
மற்றும் $\Pow{X} = Y$ என்ற தொடர்புகளை வரையறுக்கும் முழு !!{formula}s
ஐ மட்டுமே பயன்படுத்த முடியும். உண்மையில் இத்தொடர்புகள் எப்போதாவது
நிறைவுறுத்தப்படுகின்றனவா, அதாவது ஒன்றியங்களும் அடுக்குக் கணங்களும்
எப்போதும் உள்ளனவா, என்பது நமக்குத் தெரியாது. எடுத்துக்காட்டாக,
!!{sentence} $\lforall[x][\lexists[y][\Pow{x} = y]]$ என்பது~$\Log{ZFC}$
இன் மற்றோர் அடிகோள் (அடுக்குக் கண அடிகோள்).

இப்போது வரிசைச் சோடிகள் அல்லது சார்புகள் பற்றிப் பேசுவது எப்படி?
வரிசைச் சோடிகளையும் சார்புகளையும் சிறப்புவகைக் கணங்களாக எவ்வாறு
கருதலாம் என்பதை விளக்க வேண்டும். வரிசைச் சோடி $\tuple{x, y}$ ஐக்
கணம் $\{\{x\}, \{x, y\}\}$ ஆக வரையறுப்பது ஒரு வழி. ஆனால் முன்புபோல,
இக்கணத்திற்குப் பெயரிடும் !!a{function} ஒன்றை அறிமுகப்படுத்த முடியாது;
தொடர்பு $\tuple{x, y} = z$, அதாவது $\{\{x\}, \{x, y\}\} = z$ ஐ மட்டுமே
வரையறுக்க முடியும்:
\[
\lforall[u][(u \in z \liff (\lforall[v][(v \in u \liff v = x)] \lor
  \lforall[v][(v \in u \liff (v = x \lor v = y))]))]
\]
இது,~$z$ இன் !!{element}s~$u$ என்பவை $x$ ஐ மட்டுமே !!{element} ஆகக்
கொண்ட அல்லது $x$ மற்றும்~$y$ ஆகியவற்றை மட்டுமே !!{element}s ஆகக் கொண்ட
கணங்களே என்று கூறுகிறது (அதாவது, $\{x\}$ அல்லது $\{x, y\}$ உடன்
முற்றொத்த கணங்கள்). இது கிடைத்ததும், $X \times Y = Z$ போன்ற மேலதிக
கூற்றுகளைச் சொல்லலாம்:
\[
\lforall[z][(z \in Z \liff \lexists[x][\lexists[y][(x \in
        X \land y \in Y \land \tuple{x, y} = z)]])]
\]

சார்பு $f \colon X \to Y$ ஐத் தொடர்பு $f(x) = y$ ஆக, அதாவது சோடிகளின்
தொகுப்பு~$\Setabs{\tuple{x,y}}{f(x) = y}$ ஆகக் கருதலாம். அப்போது
கணம்~$f$ ஆனது $X$ இலிருந்து $Y$ க்கான சார்பு எனக் கூறுவதற்கான
நிபந்தனைகள்: (a) அது $\subseteq X \times Y$ என்ற தொடர்பு; (b) அது
முழுமையானது, அதாவது எல்லா $x \in X$ க்கும்
$\tuple{x, y} \in f$ ஆகும் ஏதேனும் $y \in Y$ உள்ளது; (c) அது
சார்புத்தன்மையுடையது, அதாவது $\tuple{x, y}, \tuple{x, y'} \in f$
எனும்போதெல்லாம் $y = y'$ (சார்பின் மதிப்புகள் தனித்திருக்க வேண்டும்).
எனவே ``$f$ ஆனது $X$ இலிருந்து $Y$ க்கான சார்பு'' என்பதைப் பின்வருமாறு
எழுதலாம்:
\begin{align*}
 \lforall[u][(u \in f \lif {}] & \lexists[x][\lexists[y][(x \in X \land y \in
      Y \land \tuple{x, y} = u)]]) \land {}\\
 \lforall[x][(x \in X \lif {}] &
      (\lexists[y][(y \in Y \land \mathrm{maps}(f, x, y))] \land {}\\
& (\lforall[y][\lforall[y'][((\mathrm{maps}(f, x, y) \land
    \mathrm{maps}(f, x, y')) \lif y = y')]]))
\end{align*}
இங்கு $\mathrm{maps}(f, x, y)$ என்பது
$\lexists[v][(v \in f \land \tuple{x, y} = v)]$ இன் சுருக்கம்
(இந்த !!{formula} ஆனது ``$f(x) = y$'' என்பதை வெளிப்படுத்துகிறது).

$f\colon X \to Y$ ஆனது !!{injective} என்பதை வெளிப்படுத்துவதும் இப்போது
கடினமல்ல; எடுத்துக்காட்டாக:
\begin{multline*}
 f \colon X \to Y \land \lforall[x][\lforall[x'][((x \in X \land x' \in
     X \land {}]] \\
 \lexists[y][(\mathrm{maps}(f, x, y) \land \mathrm{maps}(f,
        x', y))]) \lif x = x')
\end{multline*}
சார்பு~$f\colon X \to Y$ ஆனது !!{injective} iff, $f$ ஆனது
$x, x' \in X$ இரண்டையும் ஒரே~$y$ க்கு அனுப்பும்போதெல்லாம் $x = x'$.
இந்த வாய்பாட்டை $\mathrm{inj}(f, X, Y)$ என்று சுருக்கினால், கான்டரின்
தேற்றம் போன்ற சிக்கலான ஒன்றைக் கணக் கோட்பாட்டின் மொழியிலேயே கூறலாம்:
$\Pow{X}$ இலிருந்து $X$ க்கு !!{injective} சார்பு எதுவும் இல்லை:
\[
\lforall[X][\lforall[Y][(\Pow{X} = Y \lif
    \lnot\lexists[f][\mathrm{inj}(f, Y, X)])]]
\]

ஒவ்வொரு வரையறுக்கும் பண்புக்கும் ஒரு கணத்தின் இருப்பை உறுதிசெய்யும்
மற்றோர் அடிகோள் கணக் கோட்பாட்டிற்குத் தேவை என்று ஒருவர் நினைக்கலாம்.
$!A(x)$ ஆனது மாறி~$x$ கட்டற்றதாக உள்ள கணக் கோட்பாட்டு வாய்பாடு எனில்,
பின்வரும் !!{sentence} ஐக் கருதலாம்:
\[
\lexists[y][\lforall[x][(x \in y \liff !A(x))]].
\]
இந்த !!{sentence},~$!A(x)$ ஐ நிறைவுறுத்தும் அந்த~$x$ களை மட்டுமே
!!{element}s ஆகக் கொண்ட கணம்~$y$ ஒன்று உள்ளது என்று கூறுகிறது.
இத்திட்டம் ``உட்கொள்ளல் கொள்கை'' எனப்படுகிறது. இது மிகவும் பயனுள்ளதாகத்
தோன்றுகிறது; ஆனால் முரணுடையது. $!A(x) \ident \lnot x \in x$ என்று
எடுத்தால், உட்கொள்ளல் கொள்கை பின்வருமாறு கூறுகிறது:
\[
\lexists[y][\lforall[x][(x \in y \liff x \notin x)]],
\]
அதாவது, தமக்குத்தாமே !!{element}s ஆக இல்லாத எல்லாக் கணங்களின் கணம்
உள்ளது என்று கூறுகிறது. அத்தகைய கணம் இருக்க முடியாது---இதுவே ரசலின்
முரண்புதிர். உண்மையில் $\Log{ZFC}$ இக்கொள்கையின் கட்டுப்படுத்தப்பட்ட,
முரணற்ற வடிவமான பிரித்தல் கொள்கையைக் கொண்டுள்ளது:
\[
\lforall[z][\lexists[y][\lforall[x][(x \in y \liff (x \in z \land
      !A(x))]]].
\]

\begin{prob}
பின்வருவதைக் காட்டும் !!a{derivation} ஒன்றைக் கொடுத்து, உட்கொள்ளல்
கொள்கை முரணுடையது என்று காட்டுக:
\[
\lexists[y][\lforall[x][(x \in y \liff x \notin x)]] \Proves \lfalse.
\]
முதலில் $(A \lif \lnot A) \land (\lnot A \lif A)
\Proves \lfalse$ என்று காட்டுவது உதவக்கூடும்.
\end{prob}
\end{document}
